%% =====================================================================
%%  4 / 7   Results  结果                                        成员 4
%%  主要任务：呈现真实得到的结果
%%    a. 只放测到的数字（准确率/分数/耗时…），解释留给第 5 部分
%%    b. 主结果：一张表 + 一张图，正文里每张表/图都要被 \ref 引用一次
%%    c. 消融 / 敏感性 / 个案分析各放一个 \subsection
%%  篇幅建议：2--3 页；表格用 booktabs + \small，列间距 \tabcolsep 4.8--8pt
%% =====================================================================
\section{Results}
\label{sec:results}

\todo{PART 4 --- Results (Member 4): report what was actually measured, in tables and figures, without interpretation.}

% ---- 主结果表模板（换成你们的数据）----
%
% \begin{table}[t]
% \centering\small
% \setlength{\tabcolsep}{6pt}
% \begin{tabular}{@{}lccc@{}}
% \toprule
% Method & Metric A & Metric B & Metric C \\
% \midrule
% Baseline            & 00.0 & 00.0 & 00.0 \\
% \rowcolor{oursrow}[\tabcolsep][\tabcolsep]Ours & 00.0 & 00.0 & 00.0 \\
% \bottomrule
% \end{tabular}
% \caption{\textbf{One-sentence takeaway of the table.} Setting, metric
% definition, and how many runs are averaged. Best value per column in bold.}
% \label{tab:main}
% \end{table}

% ---- 建议小节 ----
%
% \subsection{Main Results}
% \label{sec:main}
% 先给结论句，再指向 Table~\ref{tab:main} 和对应的图。
%
% \subsection{Ablation Study}
% \label{sec:ablations}
% 每次只改一个因素，逐行累加时在表格里用 "+" 标记。
%
% \subsection{Case Study}
% \label{sec:cases}
% 一个成功、一个失败，各引一小段原文/日志。
